\documentclass[aps,prd,reprint,nofootinbib,longbibliography]{revtex4-2}

\usepackage{amsmath,amssymb,amsthm,bm}
\usepackage{booktabs}
\usepackage{graphicx}
\usepackage{hyperref}
\usepackage{xcolor}
\usepackage{enumitem}
\usepackage{microtype}

\newtheorem{proposition}{Proposition}
\newtheorem{definition}{Definition}
\newcommand{\M}{\mathcal{M}}
\newcommand{\X}{\mathcal{X}}
\newcommand{\F}{\mathcal{F}}
\newcommand{\dd}{\mathrm{d}}

\begin{document}

\title{Projection Folds and Apparent Pair Creation:\\
A Computational Falsification Framework for Hidden-Manifold Models}

\author{Andrew H. Bond}
\affiliation{San Jos\'e State University, San Jos\'e, California, USA}
\date{August 3, 2026}

\begin{abstract}
A smooth trajectory can intersect a constant-time slice twice after its projection
into spacetime develops a nondegenerate fold.  An observer using spacetime time as
the evolution coordinate then sees a particle-antiparticle-like pair appear, even
though the underlying trajectory neither begins nor branches.  The local statement
is exact: the two projected branches have opposite time orientation and preserve a
signed intersection number.  It is also old in spirit, lying close to
Stueckelberg's worldline interpretation of pair creation.  The open question is not
whether folds can depict pairs, but whether a local hidden-manifold dynamics can
predict their statistics without putting quantum-field-theoretic answers into the
projection, initial measure, or selection rule.

We formulate that question as a falsification campaign.  Exact monotone, fold,
degenerate, and double-fold controls validate branch counting and event location.
A conserved Hamiltonian toy model tests how easily generic local dynamics produces
coordinate-time folds and serves as a negative control for universal rate claims.
A prior-art replication anchors the campaign in Stueckelberg--Horwitz--Piron
dynamics.  The first physical gate is a held-out Schwinger-scaling challenge:
a candidate model must recover the exponent $\exp[-\pi m^2/|qE|]$, conservation,
complete outcome accounting, observer covariance, and gauge independence without
retuning.  We give an implementation plan for MATLAB on a local research server
and containerized CPU/GPU ensembles on the National Research Platform.  The
framework is designed to make a clean negative publishable: a fold may remain a
useful kinematic representation even if it fails as a mechanism of quantum pair
production.
\end{abstract}

\maketitle

\section{Introduction}

Particle creation is often drawn as a worldline that turns in coordinate time.  A
single continuous curve, parameterized by an evolution variable independent of
spacetime time, can appear on a constant-time slice as two objects with opposite
time orientation.  Stueckelberg used this construction explicitly in 1941, and
modern Stueckelberg--Horwitz--Piron (SHP) dynamics retains an invariant evolution
parameter $\tau$ independent of the spacetime coordinates
\cite{Stueckelberg1941,Land2017}.  Worldline methods also give a standard
semiclassical description of Schwinger production, where closed Euclidean
trajectories determine the nonperturbative exponent
\cite{Schwinger1951,DunneSchubert2005}.

The geometric picture invites a stronger hypothesis.  Suppose observed spacetime
is a projection of a larger state manifold.  Could apparent pair production be a
stable singularity of that projection rather than a fundamental creation event?
The question is mathematically coherent, but it is easy to answer circularly.  A
model can always be assigned a projection, hidden-state density, or stopping rule
that produces the desired number of branches.  Reproducing a picture is not the
same as deriving a rate.

This paper therefore makes a deliberately limited contribution.  It does not
propose a completed alternative to quantum field theory.  It separates:

\begin{enumerate}[leftmargin=1.5em]
\item an exact kinematic statement about projection folds;
\item a dynamical question about the flux of trajectories through a fold set; and
\item a physical question about whether that flux reproduces quantum pair
production.
\end{enumerate}

We then specify a computational campaign in which each level has its own witness
and falsification bar.  The methodology follows a simple rule: a claim may not
exceed what the recorded observables can distinguish.  In particular, every
trajectory is counted.  A model is not permitted to discard no-fold, escape,
recrossing, or numerically inconvenient outcomes and then interpret the retained
subset as a production rate.

\section{Projection folds}

Let $\M$ be a smooth state manifold and let
\begin{equation}
  z:I\rightarrow\M
\end{equation}
be a smooth trajectory parameterized by an evolution variable $\tau$.  A smooth
observation map
\begin{equation}
  \pi:\M\rightarrow\X
\end{equation}
produces an observed spacetime trajectory $x^\mu(\tau)=\pi^\mu(z(\tau))$.  An
observer or detector defines a time functional
\begin{equation}
  T_u(x)=u_\mu x^\mu,
\end{equation}
where $u$ specifies the observer's time direction.  Write
$f(\tau)=T_u(x(\tau))$.

\begin{definition}[Worldline-time fold]
A point $\tau_0$ is a nondegenerate worldline-time fold when
\begin{equation}
  f'(\tau_0)=0,
  \qquad
  f''(\tau_0)\neq 0.
\end{equation}
\end{definition}

The singularity belongs to the composition $T_u\circ\pi\circ z$.  Spacetime itself
need not be degenerate.  Nor does the fold require many hidden dimensions; one
independent evolution parameter is enough.  Higher dimension earns explanatory
weight only if it predicts the distribution or dynamics of folds.

\begin{proposition}[Local pair multiplicity]
Let $\tau_0$ be a nondegenerate worldline-time fold with
$t_0=f(\tau_0)$.  In suitable local coordinates,
\begin{equation}
  f(\tau)-t_0=\sigma(\tau-\tau_0)^2,
  \qquad \sigma\in\{+1,-1\}.
\end{equation}
For $\sigma=+1$, the number of intersections with a nearby constant-time slice
changes from zero to two as $t$ crosses $t_0$ from below.  The two regular
branches have opposite signs of $f'(\tau)$.  For $\sigma=-1$, the time-reversed
statement holds.
\end{proposition}

\begin{proof}
The one-dimensional Morse lemma gives the stated normal form.  For
$\sigma=+1$ and $t>t_0$, the two solutions are
\begin{equation}
  \tau_\pm=\tau_0\pm\sqrt{t-t_0}.
\end{equation}
Their derivatives satisfy
$f'(\tau_+)=-f'(\tau_-)\neq0$.  No real solutions exist for $t<t_0$.  The
case $\sigma=-1$ follows by reversing the direction of the time coordinate.
\end{proof}

Define the signed intersection count on a regular slice by
\begin{equation}
  N_{\mathrm{sgn}}(t)
  =\sum_{\tau_i:f(\tau_i)=t}
   \operatorname{sgn} f'(\tau_i).
\end{equation}
A simple fold changes the unsigned count by two and the signed count by zero.  If
one chooses to interpret orientation as a charge label, opposite labels follow.
That interpretation is an additional physical hypothesis; the topological result
itself establishes only opposite orientation.

The local separation of branches is also diagnostic.  If an observed spatial
coordinate has nonzero first derivative at the fold,
$x(\tau)=x_0+v(\tau-\tau_0)+\cdots$, then
\begin{equation}
  \Delta x(t)=2|v|\sqrt{|t-t_0|}+o(\sqrt{|t-t_0|}).
  \label{eq:sqrt}
\end{equation}
This square-root law is an exact instrument target.

\section{From a picture to a mechanism}

A fold permits an apparent pair event.  It does not determine when events occur,
with what probability, or with what energy and quantum numbers.  A candidate
mechanism must specify at least
\begin{equation}
  (\M,\mathcal{D},\pi,\mu_0),
\end{equation}
where $\mathcal{D}$ is a local dynamics on $\M$ and $\mu_0$ is the initial-state
measure.  The critical set is
\begin{equation}
  \F_u=\left\{z:\frac{\dd}{\dd\tau}
  T_u(\pi(z(\tau)))=0\right\},
\end{equation}
and the putative production rate is a trajectory flux through the nondegenerate
part of $\F_u$.

The phrase ``anti-hub'' is not used here.  Near a fold, a projection generally
concentrates rather than depletes observed density because the projection Jacobian
vanishes.  The mathematically relevant objects are projection multiplicity,
critical-set flux, and orientation.

\subsection{Anti-circularity conditions}

A model is excluded from claim-bearing evaluation if it contains any of the
following in explicit or equivalent form:

\begin{itemize}[leftmargin=1.5em]
\item a rule that inserts folds at a desired rate;
\item the Schwinger exponent in the initial-state weights or retention rule;
\item training and evaluation on the same field, mass, and charge region;
\item postselection on trajectories that fold or resemble particles;
\item a coordinate-time choice transformed without the observer and detector;
\item a charge or energy assignment chosen only after inspecting the branches.
\end{itemize}

A long parameter search can otherwise manufacture almost any radial or exponential
law.  The scientific object is not the best-fitting hidden geometry, but the
held-out predictions of a geometry fixed for independent reasons.

\section{Controls and candidate dynamics}

\subsection{Exact instrument controls}

The measurement stack begins with six controls.

\begin{table}[h]
\caption{Instrumentation controls.  The cubic critical point prevents a zero of
$\dd t/\dd\tau$ from being treated automatically as a pair fold.}
\begin{ruledtabular}
\begin{tabular}{lll}
Control & Time map & Required result \\
\hline
N0 & $t=\tau$ & no folds \\
P0 & $t=\tau^2$ & $0\rightarrow 2$ branches \\
P1 & $t=-\tau^2$ & $2\rightarrow 0$ branches \\
D0 & $t=\tau^3$ & degenerate, not a fold \\
M0 & $t=\tau^3-\tau$ & two critical points recovered \\
S0 & circular instanton action & $R_*=m/|qE|$, $S_*=\pi m^2/|qE|$ \\
\end{tabular}
\end{ruledtabular}
\end{table}

Control S0 uses the reduced circular worldline action
\begin{equation}
  S(R)=2\pi mR-\pi|qE|R^2.
\end{equation}
Its stationary radius and action are
\begin{equation}
  R_*=\frac{m}{|qE|},
  \qquad
  S_*=\frac{\pi m^2}{|qE|}.
\end{equation}
It verifies that the rate-analysis code can recover the standard semiclassical
scaling before it is applied to a new dynamics.

\subsection{A conserved negative control}

We use the toy Hamiltonian
\begin{equation}
\begin{split}
H={}&\frac{1}{2}(p_t^2+p_x^2+p_u^2)
+\frac{1}{2}\omega_t^2t^2
+\frac{1}{2}\omega_u^2u^2\\
&+\frac{\lambda}{4}u^4+gEtu.
\end{split}
\label{eq:toy}
\end{equation}
The variable $u$ is hidden, $\tau$ is the Hamiltonian evolution parameter, and the
coordinate-time velocity is $\dd t/\dd\tau=p_t$.  Folds occur whenever $p_t$
crosses zero with nonzero acceleration.

Equation~\eqref{eq:toy} is intentionally not relativistic and is not proposed as a
model of QED.  It asks a narrower question: how often does a generic local,
energy-conserving system produce projection folds, and how sensitive is the answer
to initial-state measure, slicing, and numerical method?  A likely result is that
folds are easy to obtain but their flux has no universal Schwinger dependence.
That negative is useful because it separates the stable geometry of a fold from the
far stronger physics of a pair-production rate.

\subsection{Prior-art replication}

Before modifying SHP dynamics, the campaign reproduces one published classical pair
trajectory and its total conservation accounting \cite{Land2017}.  The equations,
gauge-field conventions, initial conditions, and numerical tolerances are frozen in
a replication record.  Failure to reproduce the selected case blocks novelty
claims.  This stage prevents the proposed framework from quietly rediscovering a
known worldline construction under new terminology.

\section{The experiment campaign}

\subsection{PF-0: event-location audit}

PF-0 measures branch counts, signed counts, critical-point classification, and
Eq.~\eqref{eq:sqrt}.  The instrument fails if it misses a simple fold, labels the
cubic control as a fold, or changes the signed intersection number outside the
sealed numerical tolerance.

\subsection{PF-1: structural stability}

Smooth random perturbations are added to the exact fold.  The campaign measures the
persistence of a nearby nondegenerate critical point, the fitted square-root
exponent, and the condition number of event location.  Numerical precision and
sampling density are swept independently.  A positive result is not pair physics;
it is confirmation that the code sees the stable singularity it claims to see.

\subsection{PF-2: generic local dynamics}

Ensembles of Eq.~\eqref{eq:toy} are run over field-like coupling, oscillator scales,
nonlinearity, initial-state distribution, and observer slicing.  Every trajectory
receives an outcome: no fold, fold, recrossing, escape, or numerical failure.  The
primary quantities are fold flux, branch lifetime, creation--annihilation balance,
and Hamiltonian drift.  The result is rejected as a mechanism if an apparent law is
explained by the initial measure, discarded trajectories, or coordinate choice.

\subsection{PF-3: Schwinger scaling challenge}

The first physical gate uses a candidate dynamics fixed before the sweep.  One
overall time-volume conversion and, if unavoidable, one rate prefactor may be
calibrated on a training region.  Contiguous ranges of $E$, $m$, and $q$ are held
out.  The target is
\begin{equation}
  \log\Gamma=\mathrm{const}-\frac{\pi m^2}{|qE|}
  \label{eq:schwinger}
\end{equation}
within a preregistered approximation regime.

The candidate is compared with polynomial-threshold, power-law, alternative
exponential, and complexity-penalized spline models.  It fails if the held-out slope
excludes $-\pi$, if a simpler alternative predicts better, if the result is not
stable under integration refinement, or if an equivalent of Eq.~\eqref{eq:schwinger}
was embedded in the initial measure or stopping rule.

\subsection{PF-4: conservation and complete accounting}

A branch-count change is not sufficient.  The underlying model must conserve its
stated total quantities.  If the background field is dynamical, energy acquired by
the projected branches must be supplied by the field.  An orientation-based charge
rule must be declared in advance and sum to the initial charge.  Missing or deleted
trajectories are reported as primary outcomes, not ignored as implementation noise.

\subsection{PF-5: covariance and gauge audit}

A coordinate-time fold may be observer dependent.  The complete setup is therefore
transformed: state, fields, projection, and detector time direction.  A passive
coordinate change must transform event worldpoints without changing invariant
rates.  Electromagnetic implementations are repeated in gauge-equivalent potentials
with the same field tensor.  Observer-dependent particle assignments are retained
only when a physical detector model predicts the corresponding difference.

\subsection{PF-6: quantum ceiling}

Only after the preceding gates pass does the campaign test scalar--spinor
differences, multipair statistics, Pauli blocking, interference in pulsed fields,
and backreaction.  A classical fold model is expected to fail at least some of
these tests.  Such failure does not refute the fold theorem; it limits the model to
a worldline representation rather than a replacement for quantum field theory.

\section{Execution on Atlas and NRP}

The reference package contains MATLAB and Python implementations.  MATLAB on Atlas
is used first because adaptive ODE solvers can locate zero crossings through event
functions, and Parallel Computing Toolbox supports parameter sweeps with
\texttt{parfor} \cite{MatlabODE45,MatlabParfor}.  Atlas also serves as the
independent reducer for selected NRP outputs.

The National Research Platform (NRP) is used for large finite ensembles.  NRP
recommends Kubernetes Jobs for computations that run to completion and prohibits
batch jobs that merely sleep while awaiting manual use \cite{NRPJobs}.  Current
cluster policy requires resource requests to match measured use and requires equal
requests and limits for large populations of pods or jobs \cite{NRPPolicies}.  The
campaign therefore keeps its full manifest in durable storage and exposes only a
bounded live working set.  Four bounds are monitored independently:

\begin{enumerate}[leftmargin=1.5em]
\item active workers;
\item Kubernetes objects and pending pods;
\item storage and network pressure; and
\item unreduced evidence backlog.
\end{enumerate}

GPU runs begin with one device per job and record device type, CUDA runtime,
precision, and image digest.  The NRP documentation supports requesting generic or
specific GPU resources, but utilization must justify the allocation
\cite{NRPGPU}.  The supplied GPU path is a vectorized fixed-step integrator for the
toy negative control; it is not used for claim-bearing SHP or Schwinger results
until cross-validated against adaptive CPU integration.

Every trial writes an immutable JSON record containing the complete scenario,
random seed, projection, observer, numerical method, all outcome counts,
conservation residuals, resource use, code commit, container digest, and hashes of
inputs and outputs.  A reducer may summarize records but may not rewrite them.

\section{What would count as a result?}

Several outcomes would be scientifically useful.

\paragraph{Expected negative.}
Simple folds are structurally stable, generic local systems produce them, but fold
flux depends on arbitrary dynamics and initial measure and fails the held-out
Schwinger law.  The conclusion would be that projection folds are a valid
kinematic representation and not a mechanism of QED pair production.

\paragraph{Stronger negative.}
A broad model class can be shown analytically or computationally to produce only
power-law or threshold fold flux unless an exponential action or equivalent
large-deviation structure is inserted.  This could become a no-go result for naive
hidden-manifold models.

\paragraph{Unexpected positive.}
One fixed, local, covariant dynamics reproduces Eq.~\eqref{eq:schwinger}, complete
conservation, gauge independence, held-out field profiles, and later quantum
signatures.  That would justify a much stronger theory paper.  A fitted $0\to2$
branch plot would not.

\section{Relation to black-hole particle creation}

The same language may later be applied to Hawking or Unruh settings, where particle
content depends on the relation between field modes and an observer's time
coordinate.  The relevant benchmark would not be the existence of two branches,
but recovery of the exponential null-coordinate mapping, Bogoliubov coefficients,
thermal spectrum, and detector response.  We leave that extension outside the
first campaign.  A model that cannot pass the constant-field Schwinger gates has no
basis for a black-hole information claim.

\section{Discussion}

A fold supplies an economical explanation of why a projected worldline can look
like a pair with opposite orientation.  It does not explain why electrons have
their measured mass, charge, spin, statistics, and production rate.  The distinction
between depiction and derivation is the center of the proposed campaign.

The likely value of this work is methodological.  Hidden-dimensional and emergent
models are often judged by whether they can display a desired curve.  Here the
model must survive exact nulls, positive controls, complete accounting, held-out
parameter regimes, independent implementations, conservation, and covariance.  A
clean rejection leaves the geometry intact while removing an unsupported physical
interpretation.  A positive result must earn each additional word in the claim.

\section*{Data and code availability}

The accompanying source package contains the manuscript, experiment plan,
preregistration template, MATLAB controls and ODE pilots, a Python reference
implementation with unit tests, and NRP Kubernetes templates.  No claim-bearing
results are included in this draft.

\bibliography{references}

\end{document}
